\section{The two auxiliary theorems}\label{sec:inputs}

Theorems~\ref{thm:smallcut} and~\ref{thm:girth} are due to Astra. Their proofs are in the
supplement~\cite{S1,S2,S3}, which has about sixty pages. This section states what they depend on
and describes the arguments. It proves nothing.

\paragraph{Status.}
The supplement consists of research reports that Astra wrote during the work, and that it checked
with several independent passes of its own. Fable~5.2 read the report~\cite{S2}, and it checked the
proof of Theorem~\ref{thm:stability} in~\cite{S3}. Of~\cite{S1} it read the statements. It did not
check the case analysis of~\cite{S1} line by line. None of the supplement is formalised. The reports contain
exact tests of their finite identities on small graphs. A reader who wants to verify
Theorem~\ref{thm:sharp} must read the supplement.

\paragraph{No circle.}
The supplement uses Theorem~\ref{thm:forest} and Theorem~\ref{thm:cycles}. It uses neither
Theorem~\ref{thm:half} nor Theorem~\ref{thm:sharp}.

\subsection{Paths between two vertices}

For two vertices $s,t$ of a connected multigraph $N$ let $P_{st}(N)$ be the number of paths from
$s$ to $t$. By Theorem~\ref{thm:exact}(a), $P_{st}(N)\le2^{\beta(N)}$. Cycles through an edge are
paths: if $e=st$ is an edge of $G$, then $c_e(G)=P_{st}(G-e)$. So Lemma~\ref{lem:edge} reads
$c(G)=c(G-e)+P_{st}(G-e)$. The supplement studies graphs in which $P_{st}$ is more than half of its
largest possible value.

\begin{theorem}[Stability of the number of paths; supplement~\cite{S3}]\label{thm:stability}
For every $\varepsilon\in(0,\half]$ there is a constant $C_\varepsilon$ with the following
property. Let $N$ be a connected graph of rank $b$ and let $s\ne t$ be two of its vertices. If
$P_{st}(N)\ge(\half+\varepsilon)2^{b}$, then every path from $s$ to $t$ has at least
$b-C_\varepsilon$ edges.
\end{theorem}

This is a stable form of Theorem~\ref{thm:exact}(a), which is the case of equality with
$C=0$. Its proof uses Theorem~\ref{thm:forest} with a fixed error, through the localisation lemma
and a version of Lemma~\ref{lem:width} for the layers around a shortest path. The first version, in
which every path has at least $c(\varepsilon)\,b$ edges, is due to Fable~5.2. The additive form is
due to Astra. A second result of~\cite{S1} is a gap: $P_{st}(N)=2^{b}$ or
$P_{st}(N)\le\frac{15}{16}2^{b}$.

\subsection{Small cyclic cuts}

Let $G$ be a cubic multigraph with a cyclic cut $\delta X$ of size $k\le4$. Cut the $k$ edges. Each
side becomes a multigraph with $k$ vertices of degree two, the \emph{ports}. A cycle of $G$ lies
on one side, or it crosses the cut. A crossing cycle meets each side in a set of disjoint paths
between ports. So the number of crossing cycles is a sum, over the ways to pair the ports, of
products of two numbers: the number of such systems of paths on one side and on the other side.
The proof of Theorem~\ref{thm:smallcut} bounds these numbers.

For $k=2$ there is one pair of ports on each side. The crossing cycles number
$P_{st}(N_1)P_{st}(N_2)$. If both factors are at most about half of their maximum, the product is
at most about $\frac14\,2^{r-1}$, and Theorem~\ref{thm:cycles} bounds the cycles inside the two
sides. If one factor is large, Theorem~\ref{thm:stability} applies to that side. It shows that the
side is thin: its vertices lie in layers of bounded size around a shortest path. A count along the
layers then shows that such a side has only $O(2^{b/2})$ cycles. This is the
\emph{path--cycle tradeoff} of~\cite{S1}.

For $k=3$ and $k=4$ there are several pairs of ports, and for $k=4$ also systems of two paths. The
report~\cite{S1} proves inequalities between these counts on one side, classifies the sides for
which the counts are near their maxima, and evaluates the sum for each type. In the extremal type
both sides are chains of four-cycles, as in the graphs of Horak and Kim. There the largest terms of
the sum cancel exactly, and the remaining term has the order $r2^{r/2}$. The count of pairs of
disjoint cycles in a thin side is the source of the factor $r$.

\subsection{Bounded girth}

Let $G$ be a cubic multigraph without loops, of girth $g\ge5$, with no cyclic cut of size at most
four. Let $C=v_0v_1\cdots v_{g-1}$ be a shortest cycle and $e=v_0v_{g-1}$. By
Lemma~\ref{lem:edge}, $c(G)=c(G-e)+P_{v_0v_{g-1}}(G-e)$, and
$c(G-e)\le(\half+o(1))2^{r-1}$ by Theorem~\ref{thm:cycles}. The report~\cite{S2} sorts the paths
from $v_0$ to $v_{g-1}$ in $G-e$ by the vertex $v_i$ at which they leave the path
$v_0v_1\cdots v_{g-1}$. After the deletion of $v_0,\dots,v_i$, the rest of the path lies in a
multigraph of rank at most $r-2-i$. If it had more than $(\half+\varepsilon)$ of the largest
possible number of paths, a consequence of Theorem~\ref{thm:stability} would give a cyclic cut of
size at most four in $G$. So the branch at $v_i$ contributes at most about
$\half\,2^{r-2-i}$ paths. The sum over $i<g-1$ is a geometric series that falls short of
$\half\,2^{r-1}$ by $2^{r-g-1}$. This gives
\[
  c(G)\;\le\;2^{r-1}-2^{r-g-1}+o_g(2^{r}),
\]
which is Theorem~\ref{thm:girth} for each fixed girth. For girth at most four, a shortest cycle
itself gives a cyclic cut of size at most four when $r\ge9$.

The deficit $2^{r-g-1}$ decreases with the girth. So this argument cannot treat graphs whose girth
grows with $r$. Theorem~\ref{thm:half} treats them.
